\documentclass{situationist-preprint}

\title{Degrees of Moral Status}
\msid{ETH-07-c}
\msfamily{ETH-07}
\msdiscipline{Normative and Applied Ethics}
\mstype{Consequence}
\mskeywords{moral status; degrees of status; capacity for welfare; equal consideration; hierarchy}

\begin{document}
\maketitle

\begin{abstract}
ETH-07-a argues that sentience is necessary and sufficient for moral status and that the degree of a being's status varies with its capacity for welfare. We draw out what this implies about degrees. Status comes in degrees without thresholds above the minimal one; the ordering of beings by status is fixed, while its cardinal shape is fixed by equal consideration; and degrees of status do not multiply the weight of interests but summarize how much can be at stake for a being. It follows that a given harm counts the same whoever suffers it, and that differences of status bear on what we owe only through differences in what beings can gain or lose.
\end{abstract}

\section{Four consequences}

We use the framework of ETH-07-a: $\sigma(x)$ is the degree of moral status of $x$ (Definition~2.1), $c(x)$ its capacity for welfare (Definition~2.4), and an assignment of degrees is admissible if it satisfies the interest, supervenience and monotonicity principles (Definition~2.7).

\begin{proposition}[No upper thresholds]\label{p:nothreshold}
There is exactly one threshold in moral status, at sentience. Above it, status varies with capacity and changes only where capacity changes.
\end{proposition}

\begin{proof}
By Theorem~3.2 of ETH-07-a, $\sigma(x) > 0$ if and only if $x$ is sentient. By Proposition~3.3 of ETH-07-a, $\sigma$ is strictly increasing in $c$ and constant where $c$ is constant. A further threshold, a capacity level at which status jumps from one degree to a different kind of status, such as personhood, would have to be a point at which beings of equal capacity on either side differ in status, contrary to Principle~P2 of ETH-07-a, or a point at which status changes while capacity does not.
\end{proof}

Some protections apply only to beings with certain capacities, but this is a difference in what can be at stake for them, not in the kind of status they have.

\begin{proposition}[Fixed order]\label{p:order}
Whatever further principle fixes the cardinal values of status, it cannot reverse or collapse any comparison of status between two beings that differ in capacity for welfare.
\end{proposition}

\begin{proof}
Any such principle selects an admissible assignment, and by Theorem~3.4 of ETH-07-a all admissible assignments agree on the order of beings.
\end{proof}

The third consequence concerns the cardinal shape. ETH-07-a argues, in its Section~5, that equal consideration of interests (Definition~5.1) fixes the degree of status of a being as the total weight of what can be at stake for it, so that status is proportional to capacity. Hierarchical assignments such as $h(x) = c(x)^2$ are admissible (Proposition~5.2 of ETH-07-a) but count capacity twice.

\begin{proposition}[Status does not multiply interests]\label{p:nomult}
Under equal consideration, a welfare change of a given magnitude has the same moral weight whichever being undergoes it. Degrees of status therefore make a difference to what is owed to a being only through differences in the magnitude of the welfare changes that are possible for it.
\end{proposition}

\begin{proof}
The first sentence is Definition~5.1 of ETH-07-a. The degree of status is, on that definition, the total of what can be at stake for the being. A difference in status between two beings is therefore a difference in the range of welfare changes open to them, and nothing in the weight of any single change depends on whose change it is.
\end{proof}

The fourth consequence is practical. A mouse and a human being in pain of the same intensity and duration have equally weighty claims to relief, although the human has the higher status \citep{Singer1993}. The human's higher status shows itself elsewhere: in the many goods, such as long-term projects, relationships and knowledge, whose loss is possible for her and not for the mouse, and so in the greater loss that death, for example, typically imposes on her \citep{McMahan2002}. The degree of status is a summary of these differences. It is not an additional factor to be applied to them.

\section{Objection and reply}

\begin{objection}
If status does not multiply interests, it does no work. The notion of degrees is then idle, and one might as well say that all sentient beings have equal status.
\end{objection}

\begin{reply}
The notion is not idle, because it answers a question that the weight of particular interests does not: how much is owed to a being in total, over the range of ways its life can go. Two beings with equal status in that sense have equal claims across every possible situation, and a being of higher status has more that can be protected or harmed. What the notion does not do is license discounting a particular interest because of whose it is. That is the error that equal consideration excludes, and it is the error of hierarchical views such as that of \citet{Kagan2019}, not a defect of graded status.
\end{reply}

\section{Conclusion}

Moral status has one threshold, a fixed order, and a cardinal shape fixed by capacity under equal consideration. Degrees of status measure how much can be at stake for a being; they do not weight what is at stake.

\begin{thebibliography}{9}\small
\bibitem[Kagan(2019)]{Kagan2019} Kagan, S. (2019). \emph{How to Count Animals, More or Less}. Oxford University Press.
\bibitem[McMahan(2002)]{McMahan2002} McMahan, J. (2002). \emph{The Ethics of Killing: Problems at the Margins of Life}. Oxford University Press.
\bibitem[Singer(1993)]{Singer1993} Singer, P. (1993). \emph{Practical Ethics}. 2nd edn. Cambridge University Press.
\end{thebibliography}

\end{document}
